\begin{definition}[4.6.19]
\label{IV.2.4.6.19-fr}
Cuando $\mathscr F$ satisface las condiciones equivalentes de \textup{(4.6.18)}, se dice
que $\mathscr F$ es geométricamente íntegro sobre $k$.
\end{definition}

\begin{proposition}[4.6.20]
\label{IV.2.4.6.20-fr}
Sean $k$ un cuerpo, $X$ un $k$-preesquema localmente noetheriano, $\mathscr F$
un $\mathcal O_X$-módulo coherente. Sea $K$ una extensión de $k$ tal que
$X\otimes_k K$ sea localmente noetheriano. Entonces, si $\mathscr F$ es
geométricamente reducido \textup{(resp.} geométricamente íntegro) sobre $k$,
$\mathscr F\otimes_k K$ es geométricamente reducido \textup{(resp.} geométricamente
íntegro) sobre $K$.
\end{proposition}

\begin{proof}
En efecto, para toda extensión finita $K'$ de $K$,
$(X\otimes_k K)\otimes_K K'=X\otimes_k K'$
es localmente noetheriano, y se ha visto en las demostraciones de (4.6.16) y
(4.6.18) que la hipótesis sobre $\mathscr F$ implica que
$\mathscr F\otimes_k K'$ es reducido (resp. íntegro), de donde se sigue la
conclusión.
\end{proof}

\begin{proposition}[4.6.21]
\label{IV.2.4.6.21-fr}
Sean $k$ un cuerpo, $X$ e $Y$ dos $k$-preesquemas localmente noetherianos tales
que $X\times_k Y$ sea localmente noetheriano. Sean $\mathscr F$ un
$\mathcal O_X$-módulo coherente y $\mathscr G$ un $\mathcal O_Y$-módulo
coherente.
\begin{enumerate}
\item[\textup{(i)}] Si $\mathscr F$ es geométricamente reducido \textup{(resp.} geométricamente
íntegro) y $\mathscr G$ es reducido \textup{(resp.} íntegro), entonces
$\mathscr F\otimes_k\mathscr G$ es reducido \textup{(resp.} íntegro).
\item[\textup{(ii)}] Si $\mathscr F$ y $\mathscr G$ son geométricamente reducidos \textup{(resp.}
geométricamente íntegros), lo mismo ocurre con
$\mathscr F\otimes_k\mathscr G$.
\end{enumerate}
\end{proposition}

\begin{proof}
\begin{enumerate}
\item[(i)] Utilizando (3.2.3), (4.6.16) y (4.6.18), es posible limitarse al caso
en que $\operatorname{Supp}(\mathscr F)=X$,
$\operatorname{Supp}(\mathscr G)=Y$, siendo $X$ geométricamente reducido
(resp. geométricamente íntegro), $Y$ reducido (resp. íntegro), $\mathscr F$ y
$\mathscr G$ sin ciclos primos asociados inmersos, y siendo $\mathscr F_x$
isomorfo a $\mathcal O_x$ en todo punto maximal $x$ de $X$ y $\mathscr G_y$
isomorfo a $\mathcal O_y$ en todo punto maximal $y$ de $Y$. Se sabe entonces
por (4.2.3) que $\mathscr F\otimes_k\mathscr G$ no tiene ciclos primos
asociados inmersos, y que los ciclos primos asociados a
$\mathscr F\otimes_k\mathscr G$ son exactamente las componentes irreducibles
de $X\times_k Y$ (4.2.5). Es posible, pues, limitarse al caso en que $X$ e $Y$
son irreducibles~; entonces, si $X$ es geométricamente reducido e $Y$ reducido,
se sabe por (4.2.4, (ii)) que $X\times_k Y$ es reducido~; si $X$ es
geométricamente íntegro e $Y$ íntegro, se sabe que $X\times_k Y$ es además
irreducible (4.5.8, (i)), y por tanto íntegro. Finalmente, todo punto maximal
$z$ de $X\times_k Y$ está por encima de $x$ y de $y$, por lo que las hipótesis
sobre $\mathscr F$ y $\mathscr G$ implican, en virtud de (\textbf I, 9.1.12) que
$(\mathscr F\otimes_k\mathscr G)_z$ es isomorfo a $\mathcal O_z$, lo que
completa la demostración en este caso.

\item[(ii)] El razonamiento es análogo, pero aquí (después de reducir la cuestión al
caso en que $X=\operatorname{Supp}(\mathscr F)$,
$Y=\operatorname{Supp}(\mathscr G)$), $X\times_k Y$ es geométricamente reducido
(resp. geométricamente íntegro) en virtud de (4.6.5, (ii)).
\end{enumerate}
\end{proof}

Las definiciones (4.6.18) y (4.6.19) se localizan~:

\begin{definition}[4.6.22]
\label{IV.2.4.6.22-fr}
Sean $k$ un cuerpo, $X$ un $k$-preesquema localmente noetheriano, $\mathscr F$
un $\mathcal O_X$-módulo coherente. Se dice que $\mathscr F$ es
\emph{geométricamente reducido} o \emph{separable} \textup{(resp. geométricamente
puntualmente íntegro)} en un punto $x\in X$ si, para toda extensión finita
radicial \textup{(resp.} finita) $k'$ de $k$, $\mathscr F\otimes_k k'$ es reducido
\textup{(resp.} íntegro) en cada uno de los puntos $x'$ de $X\otimes_k k'$ que están
por encima de $x$ \textup{(cf. (3.2.2) y (3.2.4))}.
\end{definition}

\oldpage[IV-2]{75}

Si $\mathscr F$ es reducido en $x$, existe una vecindad abierta $U$ de $x$ tal
que $\mathscr F|U$ es reducido (3.2.2). Razonando como en (4.6.16), se ve que
decir que $\mathscr F$ es geométricamente reducido en el punto $x$ equivale a
decir que $\mathscr F$ es reducido en el punto $x$ y que el subpreesquema
cerrado $Y$ definido por el anulador $\mathscr J$ de $\mathscr F$ es
\emph{geométricamente reducido en el punto} $x$; de ello resulta que existe una vecindad
abierta $U$ de $x$ tal que $\mathscr F|U$ es \emph{geométricamente reducido} (4.6.11).
Los razonamientos de (4.6.16) y (4.6.18) muestran asimismo que, si $X$ es
localmente de tipo finito sobre $k$, decir que $\mathscr F$ es geométricamente
puntualmente íntegro en el punto $x$ equivale a decir que $\mathscr F$ es
reducido en el punto $x$ y que el subpreesquema $Y$ es \emph{geométricamente
puntualmente íntegro en el punto} $x$.

